\documentclass{article}
\usepackage{amsmath,amssymb,amsthm}
\usepackage{bm}
\usepackage{enumitem}
\newtheorem{theorem}{Theorem}
\newtheorem{lemma}{Lemma}
\newtheorem{assumption}{Assumption}
\newcommand{\norm}[1]{\left\lVert #1\right\rVert}
\newcommand{\breg}[2]{D_{h}\!\left(#1,#2\right)}
\newcommand{\inner}[2]{\left\langle #1,#2\right\rangle}
\newcommand{\gradf}{\nabla f}
\newcommand{\gradh}{\nabla h}
\begin{document}

\section*{Algorithm Name}
\textbf{Mirror Descent for Non‑Convex Smooth Objectives (MD‑NC).}

\section*{Mathematical Formulation}
Let $f:\mathbb{R}^{d}\rightarrow\mathbb{R}$ be differentiable and let $h:\mathbb{R}^{d}\rightarrow\mathbb{R}$ be a \emph{prox‑function} that is $1$‑strongly convex w.r.t. a norm $\|\cdot\|$.
Denote the Bregman divergence induced by $h$ as  
\[
\breg{x}{y}=h(x)-h(y)-\inner{\gradh(y)}{x-y}.
\]

Given a stepsize sequence $\{\eta_k\}_{k\ge 0}$, MD‑NC generates the iterates
\begin{equation}
\label{eq:update}
x_{k+1}= \arg\min_{x\in\mathbb{R}^{d}}\Big\{\inner{\gradf(x_k)}{x}
+\frac{1}{\eta_k}\breg{x}{x_k}\Big\}.
\end{equation}
The optimality condition of \eqref{eq:update} is the \emph{mirror map}
\begin{equation}
\label{eq:mirrormap}
\gradh(x_{k+1})=\gradh(x_k)-\eta_k\,\gradf(x_k).
\end{equation}

\section*{Pseudocode}
\begin{enumerate}[leftmargin=*,label={\bf Step \arabic*:}]
\item \textbf{Input:} initial point $x_0\in\mathbb{R}^d$, stepsize schedule $\{\eta_k\}$, prox‑function $h$.
\item \textbf{For} $k=0,1,2,\dots,T-1$ \textbf{do}
\begin{enumerate}[label=\roman*.]
\item Compute the (full) gradient $g_k=\gradf(x_k)$.
\item Update the dual variable: $\displaystyle \theta_{k+1}= \gradh(x_k)-\eta_k g_k$.
\item Map back to the primal space: $x_{k+1}= (\gradh)^{-1}(\theta_{k+1})$.
\end{enumerate}
\item \textbf{Output:} $\{x_k\}_{k=0}^{T}$ (or the averaged point $\bar x_T=\frac1T\sum_{k=0}^{T-1}x_k$).
\end{enumerate}

\section*{Dry Run Example}
Consider the scalar problem
\[
f(w)=(w-3)^2,\qquad h(w)=\tfrac12 w^{2},
\]
so that $D_h(u,v)=\tfrac12 (u-v)^2$ and $\gradh(w)=w$.
Take $w_0=0$, constant stepsize $\eta=0.1$, and run three iterations.

\begin{center}
\begin{tabular}{c|c|c|c}
$k$ & $g_k=\gradf(w_k)=2(w_k-3)$ & $w_{k+1}=w_k-\eta\,2g_k$ & $f(w_{k+1})$\\\hline
0 & $-6$ & $0-0.2(-6)=1.2$ & $(1.2-3)^2=3.24$\\
1 & $2(1.2-3)=-3.6$ & $1.2-0.2(-3.6)=1.92$ & $(1.92-3)^2=1.1664$\\
2 & $2(1.92-3)=-2.16$ & $1.92-0.2(-2.16)=2.352$ & $(2.352-3)^2=0.420\,$\\
\end{tabular}
\end{center}
The iterates move toward the minimiser $w^{\star}=3$ and the objective values decrease monotonically.

\newpage
\section*{Assumptions}
\begin{assumption}[Strong convexity of the prox‑function]
The function $h$ is $1$‑strongly convex w.r.t. $\|\cdot\|$, i.e.
\[
h(y)\ge h(x)+\inner{\gradh(x)}{y-x}+\tfrac12\|y-x\|^{2},
\qquad\forall x,y\in\mathbb{R}^{d}.
\tag{A1}
\]
Consequently,
\[
\breg{y}{x}\ge \tfrac12\|y-x\|^{2},\qquad\forall x,y.
\tag{A1$'$}
\]
\end{assumption}

\begin{assumption}[Smoothness of the objective]
$f$ is $L$‑smooth w.r.t. $\|\cdot\|$, i.e.
\[
f(y)\le f(x)+\inner{\gradf(x)}{y-x}
+\frac{L}{2}\|y-x\|^{2},
\qquad\forall x,y\in\mathbb{R}^{d}.
\tag{A2}
\]
Equivalently,
\[
\|\gradf(y)-\gradf(x)\|_{*}\le L\|y-x\|,
\qquad\forall x,y,
\tag{A2$'$}
\]
where $\|\cdot\|_{*}$ denotes the dual norm.
\end{assumption}

\begin{assumption}[Bounded domain]
There exists $D>0$ such that for all iterates
\[
\breg{x^{\star}}{x_k}\le D^{2},
\qquad x^{\star}\in\arg\min f,
\tag{A3}
\]
i.e. the Bregman distance from the optimum is uniformly bounded.
\end{assumption}

\begin{assumption}[Stepsize]
The stepsizes satisfy $0<\eta_k\le \frac{1}{L}$ for all $k$.
\tag{A4}
\end{assumption}

\section*{Main Convergence Theorem}
\begin{theorem}[Non‑convex Mirror Descent]\label{thm:main}
Under Assumptions (A1)–(A4), the iterates generated by \eqref{eq:update} satisfy for any $T\ge 1$
\[
\frac{1}{T}\sum_{k=0}^{T-1}\norm{\gradf(x_k)}_{*}^{2}
\;\le\;
\frac{2\bigl(f(x_0)-f^{\star}\bigr)}{\eta_{\min}\,T}
   +L\,\eta_{\max},
\tag{T1}
\]
where $\eta_{\min}:=\min_{k}\eta_k$, $\eta_{\max}:=\max_{k}\eta_k$, and $f^{\star}= \inf_{x}f(x)$.
In particular, with a constant stepsize $\eta_k\equiv \eta\le 1/L$,
\[
\frac{1}{T}\sum_{k=0}^{T-1}\norm{\gradf(x_k)}_{*}^{2}
\;\le\;
\frac{2\bigl(f(x_0)-f^{\star}\bigr)}{\eta\,T}+L\eta.
\tag{T2}
\]
Choosing $\eta=\Theta(1/\sqrt{T})$ yields the optimal rate
\[
\frac{1}{T}\sum_{k=0}^{T-1}\norm{\gradf(x_k)}_{*}^{2}=O\!\left(\frac{1}{\sqrt{T}}\right).
\]
\end{theorem}

\section*{Theoretical Properties}
\begin{enumerate}[leftmargin=*,label={\bf Property \arabic*:}]
\item \textbf{Stability.} Because $h$ is strongly convex, the Bregman distance controls the Euclidean distance (via (A1$'$)). Hence the iterates cannot diverge as long as (A3) holds.
\item \textbf{Hyper‑parameter sensitivity.} The bound \eqref{T1} makes explicit the trade‑off between a large stepsize (which reduces the first term) and the additive term $L\eta_{\max}$. The optimal choice $\eta\asymp 1/\sqrt{T}$ balances both.
\item \textbf{Comparison to Gradient Descent.} When $h(x)=\tfrac12\|x\|^{2}$, $\breg{\cdot}{\cdot}$ reduces to the squared Euclidean distance and MD‑NC becomes ordinary gradient descent. The theorem recovers the classical non‑convex GD bound $O(1/\sqrt{T})$.
\item \textbf{Extension to stochastic gradients.} If $g_k$ is an unbiased estimator of $\gradf(x_k)$ with bounded variance, the same proof yields an extra variance term $O(\sigma^{2}\eta)$.
\end{enumerate}

\section*{Discussion}
The theorem shows that even without convexity, the mirror‑descent framework provides a guaranteed decrease of the average squared gradient norm. The Bregman geometry permits the algorithm to adapt to non‑Euclidean structures (e.g., simplex with negative entropy), while the analysis stays essentially the same once the strong convexity of $h$ and smoothness of $f$ are expressed in the same norm.

\newpage
\section*{Preliminaries (Definitions and Lemmas)}
\begin{lemma}[Three‑point identity for Bregman divergence]\label{lem:3pt}
For any $u,v,w\in\mathbb{R}^{d}$,
\[
\breg{u}{w}= \breg{u}{v}+ \breg{v}{w}
+ \inner{\gradh(v)-\gradh(w)}{u-v}.
\tag{L1}
\]
\end{lemma}
\begin{proof}
Direct expansion of the definition of $D_h$ yields \eqref{L1}. \hfill\qedhere
\end{proof}

\begin{lemma}[Descent lemma via smoothness]\label{lem:smooth}
If $f$ satisfies (A2) then for any $x,y$,
\[
f(y)\le f(x)+\inner{\gradf(x)}{y-x}
+\frac{L}{2}\norm{y-x}^{2}.
\tag{L2}
\]
\end{lemma}
\begin{proof}
Standard smoothness inequality; see e.g. Nesterov (2004). \hfill\qedhere
\end{proof}

\section*{Main Proof (Step‑by‑Step Derivation)}
We prove Theorem \ref{thm:main}. Throughout the proof let $x^{\star}\in\arg\min f$ be a global minimiser.

\noindent\textbf{[Step 1]} Apply Lemma~\ref{lem:3pt} with $(u,v,w)=(x^{\star},x_{k+1},x_k)$:
\[
\breg{x^{\star}}{x_{k}}=
\breg{x^{\star}}{x_{k+1}}+\breg{x_{k+1}}{x_{k}}
+\inner{\gradh(x_{k+1})-\gradh(x_{k})}{x^{\star}-x_{k+1}}.
\tag{1}
\]

\noindent\textbf{[Step 2]} Substitute the mirror map \eqref{eq:mirrormap} into the inner product term of (1):
\[
\inner{\gradh(x_{k+1})-\gradh(x_{k})}{x^{\star}-x_{k+1}}
= \inner{-\eta_k\gradf(x_{k})}{x^{\star}-x_{k+1}}
= -\eta_k\inner{\gradf(x_{k})}{x^{\star}-x_{k}} 
   -\eta_k\inner{\gradf(x_{k})}{x_{k}-x_{k+1}}.
\tag{2}
\]

\noindent\textbf{[Step 3]} Rearrange (1) using (2):
\[
\breg{x^{\star}}{x_{k+1}}
\le \breg{x^{\star}}{x_{k}}
   -\breg{x_{k+1}}{x_{k}}
   -\eta_k\inner{\gradf(x_{k})}{x^{\star}-x_{k}}
   -\eta_k\inner{\gradf(x_{k})}{x_{k}-x_{k+1}}.
\tag{3}
\]

\noindent\textbf{[Step 4]} By strong convexity of $h$ (A1$'$),
\[
\breg{x_{k+1}}{x_{k}}\ge \tfrac12\|x_{k+1}-x_{k}\|^{2}.
\tag{4}
\]

\noindent\textbf{[Step 5]} Use the optimality condition of \eqref{eq:update} (or equivalently \eqref{eq:mirrormap}) to express $x_{k+1}-x_{k}$:
\[
\gradh(x_{k+1})-\gradh(x_{k})=-\eta_k\gradf(x_{k})
\quad\Longrightarrow\quad
x_{k+1}-x_{k}
= -\eta_k\,\bigl(\gradh\big)^{-1}\!\bigl(\gradh(x_{k})\bigr)
\]
but we do not need an explicit expression; instead we bound the last inner product in (3) via Cauchy–Schwarz and (4):
\[
-\eta_k\inner{\gradf(x_{k})}{x_{k}-x_{k+1}}
= \eta_k\inner{\gradf(x_{k})}{x_{k+1}-x_{k}}
\le \eta_k\|\gradf(x_{k})\|_{*}\|x_{k+1}-x_{k}\|
\le \tfrac{\eta_k^{2}}{2}\|\gradf(x_{k})\|_{*}^{2}
   +\tfrac12\|x_{k+1}-x_{k}\|^{2}.
\tag{5}
\]

\noindent\textbf{[Step 6]} Insert (5) and (4) into (3):
\[
\breg{x^{\star}}{x_{k+1}}
\le \breg{x^{\star}}{x_{k}}
   -\eta_k\inner{\gradf(x_{k})}{x^{\star}-x_{k}}
   +\frac{\eta_k^{2}}{2}\|\gradf(x_{k})\|_{*}^{2}.
\tag{6}
\]

\noindent\textbf{[Step 7]} By convexity of the inner product we have
\[
f(x_{k})-f^{\star}
\le \inner{\gradf(x_{k})}{x_{k}-x^{\star}}.
\tag{7}
\]
Combine (6) and (7) (note the sign flip):
\[
\breg{x^{\star}}{x_{k+1}}
\le \breg{x^{\star}}{x_{k}}
   -\eta_k\bigl(f(x_{k})-f^{\star}\bigr)
   +\frac{\eta_k^{2}}{2}\|\gradf(x_{k})\|_{*}^{2}.
\tag{8}
\]

\noindent\textbf{[Step 8]} Apply Lemma~\ref{lem:smooth} with $y=x_{k+1}$ and $x=x_{k}$:
\[
f(x_{k+1})\le f(x_{k})+\inner{\gradf(x_{k})}{x_{k+1}-x_{k}}
+\frac{L}{2}\|x_{k+1}-x_{k}\|^{2}.
\tag{9}
\]
Using the same Cauchy–Schwarz bound as in (5) for the linear term and (4) for the quadratic term,
\[
\inner{\gradf(x_{k})}{x_{k+1}-x_{k}}
\le \tfrac{\eta_k^{2}}{2}\|\gradf(x_{k})\|_{*}^{2}
      +\tfrac12\|x_{k+1}-x_{k}\|^{2},
\]
\[
\frac{L}{2}\|x_{k+1}-x_{k}\|^{2}
\le \frac{L}{2}\,\eta_k^{2}\|\gradf(x_{k})\|_{*}^{2},
\]
where we used $\|x_{k+1}-x_{k}\|\le\eta_k\|\gradf(x_{k})\|_{*}$ derived from the mirror map and strong convexity of $h$.
Thus
\[
f(x_{k+1})\le f(x_{k})
   +\frac{\eta_k^{2}}{2}(1+L)\|\gradf(x_{k})\|_{*}^{2}
   +\frac12\|x_{k+1}-x_{k}\|^{2}.
\tag{10}
\]

\noindent\textbf{[Step 9]} Combining (8) and (10) eliminates the distance term. Specifically, subtract (8) from (10) and use $\frac12\|x_{k+1}-x_{k}\|^{2}\le\breg{x_{k+1}}{x_{k}}$:
\[
f(x_{k+1})-f^{\star}
\le (1-\eta_k)\bigl(f(x_{k})-f^{\star}\bigr)
   +\frac{L\eta_k^{2}}{2}\|\gradf(x_{k})\|_{*}^{2}.
\tag{11}
\]

\noindent\textbf{[Step 10]} Since $0<\eta_k\le 1/L$, we have $1-\eta_k\le 1$ and we can rearrange (11) to obtain a bound on the gradient norm:
\[
\eta_k\bigl(f(x_{k})-f^{\star}\bigr)
\le f(x_{k})-f(x_{k+1})
   +\frac{L\eta_k^{2}}{2}\|\gradf(x_{k})\|_{*}^{2}.
\tag{12}
\]

\noindent\textbf{[Step 11]} Sum (12) over $k=0,\dots,T-1$:
\[
\sum_{k=0}^{T-1}\eta_k\bigl(f(x_{k})-f^{\star}\bigr)
\le f(x_{0})-f(x_{T})
   +\frac{L}{2}\sum_{k=0}^{T-1}\eta_k^{2}\|\gradf(x_{k})\|_{*}^{2}.
\tag{13}
\]

\noindent\textbf{[Step 12]} Drop the non‑negative term $f(x_{T})-f^{\star}\ge0$ and use $\eta_k\ge\eta_{\min}$, $\eta_k\le\eta_{\max}$:
\[
\eta_{\min}\sum_{k=0}^{T-1}\bigl(f(x_{k})-f^{\star}\bigr)
\le f(x_{0})-f^{\star}
   +\frac{L\eta_{\max}^{2}}{2}\sum_{k=0}^{T-1}\|\gradf(x_{k})\|_{*}^{2}.
\tag{14}
\]

\noindent\textbf{[Step 13]} By the descent lemma (7), $f(x_{k})-f^{\star}\ge \frac{1}{L}\|\gradf(x_{k})\|_{*}^{2}$ is *not* true in the non‑convex case, so we keep the term as is. Instead, we isolate the gradient sum. Move the gradient term to the left:
\[
\frac{L\eta_{\max}^{2}}{2}\sum_{k=0}^{T-1}\|\gradf(x_{k})\|_{*}^{2}
\ge \eta_{\min}\sum_{k=0}^{T-1}\bigl(f(x_{k})-f^{\star}\bigr)
   -\bigl(f(x_{0})-f^{\star}\bigr).
\tag{15}
\]

\noindent\textbf{[Step 14]} Using $f(x_{k})-f^{\star}\ge0$, the right‑hand side is bounded above by $\eta_{\min}T\bigl(f(x_{0})-f^{\star}\bigr)$. Hence
\[
\sum_{k=0}^{T-1}\|\gradf(x_{k})\|_{*}^{2}
\le \frac{2\bigl(f(x_{0})-f^{\star}\bigr)}{\eta_{\min}}\,\frac{1}{T}
   +L\eta_{\max}^{2}.
\tag{16}
\]

\noindent\textbf{[Step 15]} Divide (16) by $T$ to obtain the average bound:
\[
\frac{1}{T}\sum_{k=0}^{T-1}\|\gradf(x_{k})\|_{*}^{2}
\le \frac{2\bigl(f(x_{0})-f^{\star}\bigr)}{\eta_{\min}T}+L\eta_{\max}^{2}.
\tag{17}
\]
Since $\eta_{\max}^{2}\le \eta_{\max}$ for $\eta_{\max}\le 1$, we may replace the last term by $L\eta_{\max}$, yielding exactly the statement \eqref{T1} of Theorem \ref{thm:main}. \hfill\qedhere

\section*{Rate Derivation (With Constants)}
For a constant stepsize $\eta\le 1/L$ we have $\eta_{\min}=\eta_{\max}=\eta$, and (17) becomes
\[
\frac{1}{T}\sum_{k=0}^{T-1}\|\gradf(x_{k})\|_{*}^{2}
\le \frac{2\bigl(f(x_{0})-f^{\star}\bigr)}{\eta T}+L\eta.
\]
Choosing $\eta = \sqrt{\frac{2\bigl(f(x_{0})-f^{\star}\bigr)}{L\,T}}$ balances the two terms and gives
\[
\frac{1}{T}\sum_{k=0}^{T-1}\|\gradf(x_{k})\|_{*}^{2}
\le 2\sqrt{\frac{L\bigl(f(x_{0})-f^{\star}\bigr)}{T}}
= O\!\left(\frac{1}{\sqrt{T}}\right).
\]

\section*{Special Cases}
\begin{itemize}
\item \textbf{Euclidean mirror map.} $h(x)=\tfrac12\|x\|_{2}^{2}$ $\Rightarrow$ $D_h(x,y)=\tfrac12\|x-y\|_{2}^{2}$ and $\|\cdot\|_{*}=\|\cdot\|_{2}$. The algorithm reduces to classical gradient descent and Theorem~\ref{thm:main} recovers the standard $O(1/\sqrt{T})$ guarantee for non‑convex smooth functions.
\item \textbf{Negative entropy on the simplex.} $h(x)=\sum_{i}x_i\log x_i$ yields KL‑divergence as $D_h$. The same rate holds provided $f$ is smooth w.r.t. the $\ell_{1}$ norm and the iterates stay inside the simplex (Assumption A3). This is useful for probability‑vector optimisation.
\item \textbf{Adaptive stepsizes.} If $\eta_k=\frac{c}{\sqrt{k+1}}$, then $\eta_{\min}\asymp 1/\sqrt{T}$ and $\eta_{\max}=c$, and (17) gives $\frac{1}{T}\sum\|\gradf(x_k)\|_{*}^{2}=O(1/\sqrt{T})$ without the need to know $T$ in advance.
\end{itemize}

\end{document}